\FloatBarrier
\section{A bias-aware criterion for finite measurements}\label{app:bias-aware-design}

A measurement can reduce sampling variance while increasing approximation error.
This appendix makes that tradeoff explicit for full-rank calibration and small
effect maps. It specializes classical bias-aware design~\cite{box1959design}
using the intervention error in Lemma~\ref{lem:measurement-reduction}. The
estimator, target, noise model, and measurement costs remain fixed when designs
are compared. Full column rank distinguishes this setting from the
underdetermined sparse fits elsewhere in the paper.

\subsection{From measurement error to estimation risk}

Fix a design $S$ before observing its noise. Write
\begin{equation}\label{eq:bd-model}
 y=A\theta+b+\eta,\qquad \E\eta=0,\qquad
 \E[\eta\eta^\top]=\Sigma\succ0,
\end{equation}
where $A\in\R^{m\times d}$ has full column rank, $\theta$ is the target,
$b$ is deterministic approximation or systematic error, and $\eta$ is centered
sampling noise. Correlated measurements are allowed. Let $W\succeq0$ specify
the same target loss for every design, with $\|z\|_W^2=z^\top Wz$. Define
\begin{equation}\label{eq:bd-gls}
 M=A^\top\Sigma^{-1}A,\qquad
 G=M^{-1}A^\top\Sigma^{-1},\qquad \widehat\theta=Gy.
\end{equation}
Here $\widehat\theta$ is the generalized least-squares (GLS) estimator.
The identity below uses the true covariance; a fitted covariance would require
accounting for its estimation error.

\begin{proposition}[Fixed-design risk]\label{prop:bd-risk}
Under \eqref{eq:bd-model}--\eqref{eq:bd-gls},
\begin{equation}\label{eq:bd-risk}
 \E\|\widehat\theta-\theta\|_W^2
 =\operatorname{tr}(WM^{-1})+\|Gb\|_W^2.
\end{equation}
\end{proposition}
\begin{proof}
Full column rank gives $GA=I$. The error is $Gb+G\eta$, its covariance is
$G\Sigma G^\top=M^{-1}$, and the cross term has expectation zero. Expanding
the squared norm proves the result without a Gaussian assumption.
\end{proof}

With no bias, minimizing this risk is weighted A-optimal design. D-optimality
instead maximizes $\det M$ and generally optimizes a different objective.
Under a correctly specified Gaussian model with parameter-independent covariance,
$M$ is the Fisher information. Full rank makes GLS well-defined, but unknown
bounded bias can still prevent exact identification of $\theta$.

For MT, fix the context distribution, clean states $h_c$, and accessible
directions $D_c$. For mask $a_j$ and scale $\eps_j>0$, set
\begin{equation}\label{eq:bd-target}
 \theta=\E_c[D_c^\top\nabla f_c(h_c)],\qquad
 \kappa_j^2=\E_c\|D_ca_j\|_2^2.
\end{equation}
Under the lemma's activation-space Hessian bound $L$, the mean normalized
response satisfies
\begin{equation}\label{eq:bd-taylor}
 \begin{aligned}
 F_j&=\frac{\E_c[f_c(h_c+\eps_jD_ca_j)-f_c(h_c)]}{\eps_j}\\
    &=a_j^\top\theta+r_j,\qquad
 |r_j|\leq\frac{L\eps_j}{2}\kappa_j^2.
 \end{aligned}
\end{equation}
The bound follows by integrating the quadratic Taylor remainder along each
segment and then averaging. If an additional systematic error is bounded by
$\bar\tau_j/\eps_j+\bar\nu_j$, the total envelope is
\begin{equation}\label{eq:bd-envelope}
 |b_j|\leq\beta_j
 :=\frac{L\eps_j}{2}\kappa_j^2+\frac{\bar\tau_j}{\eps_j}+\bar\nu_j.
\end{equation}
The two kinds of error enter the model differently. Deterministic error, such
as a systematic approximation in the response numerator, is part of the bias
$b_j$ and is bounded through $\bar\tau_j$. Centered sampling variation,
including the covariance created when measurements share a random baseline,
belongs in $\Sigma$. Each error is assigned to exactly one of the two terms.
Lemma~\ref{lem:measurement-reduction} bounds their magnitudes, so the zero-mean
and covariance structure in \eqref{eq:bd-model} is an additional assumption.

\begin{corollary}[Bias-aware design criterion]\label{cor:bd-design}
Let $B=\operatorname{diag}(\beta_1,\ldots,\beta_m)$ and let $g_j$ denote
column $j$ of $G$. The risk of the physical response is at most
\begin{align}
 J_{\rm box}(S)&=\operatorname{tr}(WM^{-1})+
 \max_{s\in\{-1,1\}^m}s^\top BG^\top WGBs,\label{eq:bd-box}\\
 J_{\rm box}(S)&\leq J_{\rm sum}(S)
 :=\operatorname{tr}(WM^{-1})+
 \left(\sum_j\beta_j\sqrt{g_j^\top Wg_j}\right)^2.\label{eq:bd-sum}
\end{align}
The first expression is the exact worst-case GLS risk over the rectangle
$|b_j|\leq\beta_j$. For a scalar target $v^\top\theta$, so $W=vv^\top$,
the two expressions coincide:
\begin{equation}\label{eq:bd-scalar}
 J_{\rm box}=J_{\rm sum}
 =v^\top M^{-1}v+\left(\sum_j\beta_j|v^\top g_j|\right)^2.
\end{equation}
\end{corollary}
\begin{proof}
The squared bias is a convex quadratic in $b$. Every point of the rectangle
is a convex combination of its vertices $Bs$, so its maximum occurs at a
vertex. The triangle inequality gives
$\|W^{1/2}Gb\|_2\leq\sum_j\beta_j\|W^{1/2}g_j\|_2$.
For a scalar target, choose each sign to agree with $v^\top g_j$; this attains
the bound and proves \eqref{eq:bd-scalar}.
\end{proof}

Minimizing $J_{\rm box}$ over a finite set of feasible designs therefore
selects, among the declared GLS designs, the one with the smallest worst-case
risk over the rectangle. This comparison holds the estimator family and error
model fixed; it does not optimize over all estimators or physical response
classes. For a scalar target, the criterion reduces to the closed-form sum in
\eqref{eq:bd-scalar}. For general $W$, evaluating $J_{\rm box}$ can require a
search over the $2^m$ sign vectors. The cheaper bound $J_{\rm sum}$ can
overstate the worst squared bias by up to a factor of $m$. That factor follows
from Cauchy--Schwarz, since the average over vertices equals the sum of squared
weighted-column norms; equal-norm orthogonal weighted columns attain it.

The rectangle can also be conservative. When probes lie on the same ray at
different scales, their errors arise from the same curvature and need not
attain their separate extremes jointly. Under a Hessian bound alone, however,
the rectangle can be sharp as a supremum, as the signed example below shows.
In every case, the curvature bound and covariance are inputs that must be
known or estimated, and estimating them counts toward the measurement budget.

\subsection{Why the smoothness assumption matters}

Shared curvature can make the coordinatewise envelope conservative. Suppose
the response in mask coordinates is exactly quadratic:
\begin{equation}\label{eq:bd-quadratic}
 f(z)-f(0)=\theta^\top z+\tfrac12z^\top Hz,
 \qquad H=H^\top,\quad \|H\|_{\rm op}\leq\Lambda.
\end{equation}
The constant $\Lambda$ bounds curvature in mask coordinates and differs from
the activation-space bound $L$: for a fixed direction matrix $D$, the
pulled-back Hessian $D^\top H_{\rm act}D$ has operator norm at most
$L\|D\|_{\rm op}^2$. Assuming a shared quadratic is stronger than assuming a
Hessian bound. Under that assumption, with no other systematic error and a
fixed covariance, define
\begin{equation}
 K_S=\tfrac12\sum_j\eps_j(v^\top g_j)a_ja_j^\top.
\end{equation}
Then the exact worst-case scalar risk for the declared GLS estimator is
\begin{equation}\label{eq:bd-shared}
 v^\top M^{-1}v+\Lambda^2\|K_S\|_*^2,
\end{equation}
where $\|K_S\|_*$ is the sum of absolute eigenvalues. To prove this, use
$b_j=(\eps_j/2)a_j^\top Ha_j$ to write the scalar bias as
$\operatorname{tr}(HK_S)$. If $K_S=U\operatorname{diag}(\lambda_i)U^\top$,
its absolute maximum is $\Lambda\sum_i|\lambda_i|$, attained by
$H=\Lambda U\operatorname{diag}(\operatorname{sign}\lambda_i)U^\top$.
Unlike the rectangle, this expression allows cancellation before taking
absolute values.

For a general scalar $f$ with $|f''|\leq L$, however, the central difference
has worst-case bias $Lh/2$, the same as the forward difference. Separate
Taylor bounds at $h$ and $-h$ give the upper bound. To approach it, choose
$f_\delta''(t)=L\tanh(t/\delta)$, $f_\delta(0)=0$, and
$f_\delta'(0)=\theta$. Direct integration gives
\begin{equation}\label{eq:bd-sharpness}
 \frac{f_\delta(h)-f_\delta(-h)}{2h}-\theta
 =\frac{L\delta}{h}\int_0^h\log\cosh(t/\delta)\,dt
 \longrightarrow\frac{Lh}{2}.
\end{equation}
Indeed, the bias lies between $L(h/2-\delta\log2)$ and $Lh/2$ because
$x-\log2\leq\log\cosh x\leq x$ for $x\geq0$. With $L=h=1$ and
$\delta=0.001$, the bias is approximately $0.4993073$. The function is
infinitely differentiable; smoothness without a quantitative higher-order
bound does not guarantee cancellation.

Central differences improve on this worst-case rate under a stronger
smoothness assumption. If the activation-space Hessian is Lipschitz with
constant $L_3$ on the signed segments, the standard remainder bound gives
\begin{equation}\label{eq:bd-central}
 \left|\frac{f(h_0+\eps u)-f(h_0-\eps u)}{2\eps}
       -\nabla f(h_0)^\top u\right|
 \leq\frac{L_3\eps^2}{6}\|u\|_2^3.
\end{equation}
The second-order terms of the two Taylor expansions cancel, and each
third-order remainder is at most $L_3\eps^3\|u\|_2^3/6$. Subtracting the
expansions and dividing by $2\eps$ gives the bound, which a one-dimensional
cubic attains. A shared quadratic therefore gives zero bias, while a
Lipschitz Hessian gives bias of order $\eps^2$~\cite{shi2022finite}.
The scheme needs both signs of the intervention to be feasible, which may
fail for patching or ablation, and both evaluations count toward its cost.
Lemma~\ref{lem:measurement-reduction} needs only the Hessian bound; the
higher-order rate needs the additional smoothness stated here, which the
clipped HMM readout of Sec.~\ref{sec:aggregate} lacks on paths that cross its
clipping boundaries.

\subsection{An equal-budget comparison}

Consider $f(t)=\theta t+(q/2)t^2$ with $|q|\leq\Lambda$, a known baseline
$f(0)=0$, and a budget of four evaluations, each with independent noise
variance $\sigma^2$. If all four are taken at the same positive step $h$,
the slope estimate has variance $\sigma^2/(4h^2)$ and worst-case bias
$\Lambda h/2$, so its worst-case risk is
$R(h)=\sigma^2/(4h^2)+\Lambda^2h^2/4$.
At $\Lambda=1$, $\sigma=0.1$, and $h\in[0.1,1]$, information-only selection
chooses $h=1$ with risk $0.2525$, while minimizing $R$ gives $h=\sqrt{0.1}$
and risk $0.05$. Two evaluations at each sign instead cancel the shared
quadratic bias and give risk $0.0025$ at $h=1$.
The symmetric control therefore outperforms the best one-sided design here;
the cancellation requires the shared quadratic, not merely a Hessian bound.
If the baseline evaluation were itself noisy, it would add cost and
covariance terms to each design.

\subsection{Finite-response calibration has a different target}

The HMM calibration in Sec.~\ref{sec:step0} estimates a gain for finite
responses, not the local derivative. Fix its scale and evaluation distribution
$Q$ before choosing measurements. Let $p(a)$ be the known scalar prediction,
$F(a)$ the mean finite response, and assume finite second moments and
$\E_Qp^2>0$. Define
\begin{equation}\label{eq:bd-gain}
 g_Q=\frac{\E_Q[pF]}{\E_Q[p^2]},\qquad
 e(a)=F(a)-g_Qp(a),\qquad \E_Q[pe]=0.
\end{equation}
For a selected set $S$, let $I=p_S^\top\Sigma_S^{-1}p_S>0$ and
$\widehat g=I^{-1}p_S^\top\Sigma_S^{-1}y_S$. Proposition~\ref{prop:bd-risk}
and the projection orthogonality give
\begin{align}
 \E(\widehat g-g_Q)^2
 &=I^{-1}+(I^{-1}p_S^\top\Sigma_S^{-1}e_S)^2,\label{eq:bd-gainrisk}\\
 \E_\eta\E_Q[(p\widehat g-F)^2]
 &=\E_Q[p^2]\E(\widehat g-g_Q)^2+\E_Q[e^2].\label{eq:bd-prediction}
\end{align}
The cross term vanishes because $\E_Q[pe]=0$ and $Q$ is fixed independently
of the fitting noise. For a gain target other than $g_Q$, the cross term
generally does not vanish.

\begin{corollary}[Raw-response envelope for gain calibration]\label{cor:bd-gain}
If $F=p+r$, $|r(a)|\leq u(a)$, and $\E_Q[|p|u]<\infty$, set
$\alpha_Q=\E_Q[|p|u]/\E_Q[p^2]$. Then
\begin{equation}\label{eq:bd-gainbound}
 |g_Q-1|\leq\alpha_Q,\qquad
 |e(a)|\leq u(a)+\alpha_Q|p(a)|.
\end{equation}
\end{corollary}
\begin{proof}
Substitution in \eqref{eq:bd-gain} gives
$g_Q-1=\E_Q[pr]/\E_Q[p^2]$. The triangle inequality bounds this by
$\alpha_Q$. Apply it again to $e=r-(g_Q-1)p$.
\end{proof}

For an exact raw attribution prediction $p(a)=\eps a^\top\theta$, Taylor's
theorem gives
\begin{equation}
 u(a)=\frac{L\eps^2}{2}\E_c\|D_ca\|_2^2.
\end{equation}
The bound on this unnormalized response scales as $\eps^2$; if the
attribution map is itself approximate, its error must be added. Applying the
corollary requires a curvature bound and the $Q$-moments of $p$, and any
finite-response measurements used to obtain them count toward the budget.
When $L$ is unknown or very conservative, the criterion becomes expensive
to apply or too loose to guide the design.

The residual envelope need not be tight even within a constant factor. If
$F=2p$ and $u=|p|$, then $g_Q=2$ and $e=0$, but the envelope is $2|p|$.
With independent equal-variance noise, selected predictors of magnitude $P$
and a constant raw envelope $U$ give the scalar bound
\begin{equation}
 \frac{\sigma^2}{mP^2}+\left(\frac{U}{P}+\alpha_Q\right)^2.
\end{equation}
For fixed $U$ and $\alpha_Q$, both terms decrease with $P$. Thus normalized
masks and a uniform envelope do not establish that large attribution responses
are poor calibration choices. Showing that saturation
makes large responses poor calibration choices would require additional
structure in the residual or direct empirical evidence.

\paragraph{Scope of implementation.}
The analysis extends to other fixed linear estimators. For an estimator
$G_0$, the variance becomes $\operatorname{tr}(WG_0\Sigma G_0^\top)$ and the
squared bias becomes $\|(G_0A-I)\theta+G_0b\|_W^2$; for ridge regression,
this includes the shrinkage bias that a direct application of the GLS result
would omit. Applying the analysis after earlier data have been used to choose
the design requires the noise assumptions to hold conditionally on that choice.
Every pilot, covariance, and curvature measurement also counts toward the cost.
Because the argument fixes the design before the fitting measurements, it
gives no guarantee for sequential selection. The results so far are analytic;
the retained-data comparison below shows how information and prediction
error can diverge.

\paragraph{A two-stage acquisition workflow.}\label{par:bd-workflow}
We outline how the criterion could be applied to a small calibration model.
The target, intervention scale, and candidate family are fixed first, and a
separate validation budget is reserved. A pilot stage of $k_1$ uniformly drawn
probes measures the finite-response mismatch, while the design matrix shows
its conditioning. The pilot reveals where the local approximation fails but
does not by itself bound curvature at unmeasured interventions.

When justified bias envelopes and a specified covariance model are available,
a refinement stage chooses a full-rank design of $k_2$ probes by minimizing
$J_{\rm box}$ or its tractable upper bound $J_{\rm sum}$ over feasible designs.
The reported selection method distinguishes an exact minimum from an
approximation. If pilot estimates replace established envelopes or covariance
assumptions, the rule is a heuristic whose sensitivity needs to be assessed.
The fixed-design analysis applies conditionally on the pilot when the
refinement design is then fixed and fresh fitting responses satisfy its
noise assumptions.

Selecting $k_2$ distinct probes from $N$ candidates allows $\binom{N}{k_2}$
subsets before constraints are applied. Using $J_{\rm sum}$ avoids the inner
sign enumeration but leaves this outer search. Greedy additions from a
full-rank seed or rank-preserving probe exchanges avoid exhaustive search,
although neither has an optimality guarantee here. Selection runtime is
therefore reported separately from measurement cost.

The refinement responses fit the model, and the reserved validation
measurements assess its predictions at the intended scale. A cost-matched
comparison with uniform selection includes pilot, refinement, and validation
work, together with any measurements used to estimate covariance or bias.
We propose this workflow but have not evaluated it; the replay below tests
only fixed selection rules.

\subsection{Calibration-probe selection on retained HMM measurements}\label{app:selector-replay}

\begin{table}[htbp]
\centering
\setlength{\tabcolsep}{3pt}
\begin{tabular}{clrrr}
\toprule
Scale & Selection & \shortstack{Squared\\bias} & \shortstack{Subset\\variance} & \shortstack{Total\\MSE} \\
\midrule
$5$ & Largest-first & $0.12699$ & $0.00000$ & $0.27823$ \\
$5$ & Random        & $0.00807$ & $0.04454$ & $0.20385$ \\
$5$ & Spread        & $0.07358$ & $0.00000$ & $0.22483$ \\
\midrule
$8$ & Largest-first & $0.17359$ & $0.00000$ & $0.47405$ \\
$8$ & Random        & $0.02214$ & $0.08331$ & $0.40591$ \\
$8$ & Spread        & $0.09887$ & $0.00000$ & $0.39933$ \\
\bottomrule
\end{tabular}
\caption{Eight-probe calibration on retained HMM measurements. Entries average
the cellwise terms in Eq.~\eqref{eq:bd-empirical-decomposition}. Total MSE also
includes the best-gain error $L_Q^\star$, averaging $0.15125$ at $\eps=5$ and
$0.30046$ at $\eps=8$. Largest-first and spread are deterministic,
so their subset variance within each cell is zero; their results still vary
across cells.}
\label{tab:selector-decomposition}
\end{table}

Under a scalar linear model with independent, zero-mean, equal-variance errors,
choosing the largest
absolute predictions maximizes information and is both A- and D-optimal. Finite
responses need not follow that error model. We compare this rule with uniform
random selection and a simple alternative that spreads probes across prediction
magnitudes, using the retained HMM measurements from Sec.~\ref{sec:step0}.

All three rules see only candidate identifiers and attribution predictions.
They select from the same 96-mask pool, fit the same no-intercept gain, and are
scored on the same 32 held-out masks. The spread rule chooses unused masks nearest
to equally spaced targets over the range of absolute predictions. Random
selection uses 50 uniform draws in each cell; the same draws are used at both
scales, and smaller subsets are nested within larger ones. The primary comparison, fixed before the replay, is mean held-out MSE
at $\eps=8$ with eight probes, averaged equally across the nine cells. A probe
is one mask response averaged over 256 sequences; held-out evaluation uses 32
additional mask responses. This replay makes no new model calls.

\paragraph{Selection changes calibration error.}
At eight probes, largest-first has $16.8\%$ higher mean MSE than random at
$\eps=8$ and $36.5\%$ higher MSE at $\eps=5$
(Table~\ref{tab:selector-decomposition}). It is worse in five of nine cells at
$\eps=8$ and eight of nine at $\eps=5$. Spread improves on largest-first at both
scales, but is $10.3\%$ worse than random at $\eps=5$ and only $1.6\%$ better at
$\eps=8$. The latter MSE difference, $0.00659$, is smaller than the conditional
Monte Carlo standard error of the random mean, $0.00789$.
With four probes, largest-first beats random selection only at $\eps=8$;
with sixteen it loses at both scales, and with the full pool all three rules
select the same masks. Because this comparison redraws the random subsets,
its eight-probe means ($R^2=0.953$ at $\eps=5$ and $0.949$ at $\eps=8$)
differ slightly from those in Sec.~\ref{sec:step0} ($0.951$ and $0.948$),
which use the original draws.

\paragraph{An exact bias--variance breakdown.}
To explain these errors after the comparison, let $F$ denote the retained
response and $Q$ the uniform distribution over one cell's held-out masks.
Let $g_Q$ be its least-squares gain and $w=\E_Q[p^2]$.
Writing $L_Q(g)=\E_Q[(gp-F)^2]$ and
$L_Q^\star=L_Q(g_Q)$, projection orthogonality gives the exact identity
\begin{equation}\label{eq:bd-empirical-decomposition}
 \E_S L_Q(\widehat g_S)
 =L_Q^\star+w\bigl(\E_S\widehat g_S-g_Q\bigr)^2
   +w\operatorname{Var}_S(\widehat g_S).
\end{equation}
Here $S$ ranges over the saved subsets, with equal weight on each of the 50
random draws. The variance measures variation across these subsets, not fresh
measurement noise. The held-out gain is used only for this post-hoc diagnosis;
it is never available to a selector.

At $\eps=8$, largest-first fits a mean gain of $0.642$, compared with $0.742$
for the held-out best gain; at $\eps=5$ the corresponding values are $0.749$
and $0.886$. Random selection has less squared bias but nonzero subset variance.
Its squared-bias term accounts for about $21\%$ of its excess error above the
best-gain error at $\eps=8$. The comparison shows how calibration bias can
outweigh a variance-based selection advantage. The decomposition attributes
the excess error to bias in the fitted gain. Identifying the mechanism behind
that bias, such as saturation, and testing the bias-aware criterion
prospectively would require further measurements of curvature and covariance.
The comparison covers nine combinations of evaluation batch and mask design
from a single trained model.
\FloatBarrier
